% =====================================================================
%  24.4  Divide and conquer on trees beyond path sums -- an ordered
%  structure inside the centroid bucket.
%  Exports: sec:c18-beyond  def:c18-beyond  lst:c18-beyond
%           prob:c18-833d  lst:c18-833d
%  Solved CF 833D "Red-Black Cobweb" (2800), live-verified via the problemset
%  API (tags data structures / divide and conquer / implementation / trees).
%  Statement: tree, each edge a colour (0 red / 1 black) and a clamminess x.
%  A pair's (unique) path is kept when #(red) and #(black) "differ at most
%  twice", i.e. max(r,b) <= 2*min(r,b).  Jelliness = product over kept
%  unordered pairs of (product of x on the path), mod 1e9+7; empty -> 1.
%  KEY reduction (verified): with r=#red, b=#black on a path, the filter
%  r<=2b and b<=2r is TWO additive features u=2r-b>=0 and v=2b-r>=0; along a
%  centroid's two arms these sum, so a kept pair is a 2-D dominance pair --
%  the same shape as 293E, but the aggregate is a PRODUCT, folded as
%  prod_x w_x^{#partners(x)} and combined by G(all) / prod_arm G(arm).
%  Official sample 1 replayed by hand at centroid 2 -> 1265625 (pairs (1,3),
%  (1,5),(3,4),(4,5) with products 45,45,25,25).  Sample 2 -> 452841614.
%  Back-cites Ch.18 Fenwick (c14); the BIT is used, never re-taught.
% =====================================================================
\section{Beyond Path Sums: A Structure in the Bucket}\label{sec:c18-beyond}

\begin{previouslybox}[\(\secref{c18-count}\) and \(\secref{c18-index}\) both counted; a path \emph{product} is new]
Two sections back, the centroid bucket held a small integer array (a distance
histogram, \(\coderef{c18-count}\)); one section back it held a sorted table and a
Fenwick (a two-sided budget, \(\coderef{c18-293e}\)). Both \emph{counted} qualifying
pairs. Chapter~18 (\chref{c14}) owns the Fenwick itself; here we only bolt one into a
bucket. The step up in this section is the \emph{aggregate}: some questions are not
"how many" but "combine a per-path value over all qualifying paths" --- a product, a
sum of products, a maximum --- and the fold is not additive. CF~833D asks for a
\emph{product} over paths of a \emph{product} along each path, so we need the general
recipe: reduce a path property to per-node \emph{features} that sum across the two arms
at a centroid, decide acceptance with a dominance predicate on those features, and carry
whatever monoid the aggregate needs in the bucket.
\end{previouslybox}

\begin{topicintro}
A centroid decomposition is a divide-and-conquer, and like any one it is only as
general as the three choices you plug into it: the \emph{feature} each node contributes,
the \emph{acceptance test} for a straddling pair, and the \emph{accumulator} folded over
accepted pairs. The first two sections fixed all three (feature = distance, test = a
threshold, accumulator = count). Here we free them. The recurring trick: pick per-node
quantities that are \emph{additive along the path}, so that at the top centroid \(c\) a
pair \((x,y)\) has path-value \(f(x)+f(y)\); acceptance then says "\(\ge 0\) in two
directions," a 2-D dominance query answered by sorting one direction and sweeping a
Fenwick over the other --- the identical skeleton as \(\coderef{c18-293e}\), now with the
accumulator promoted from a count to a modular product. The same-arm correction survives
too, because a product over a set difference of pairs is a quotient of the two products
(\(\defref{c18-beyond}\)).

\textbf{Prerequisites.} \(\secref{c18-count}\) (the top-centroid claim and same-arm
subtraction), \(\secref{c18-index}\) (the 2-D dominance sweep of \(\coderef{c18-293e}\)),
and the Fenwick from Chapter~18 (\chref{c14}). Multiplicative inverses modulo a prime
(Fermat) are used to "subtract" a product; that is arithmetic, not a new data structure.

\textbf{After this section you can}
\begin{itemize}[nosep]
  \item state the general centroid D\&C template and name its three plug-ins (\(\defref{c18-beyond}\)); \hfill\lv{5}
  \item turn "the red and black counts differ by a factor of at most two" into two additive
        dominance features; \hfill\lv{5}--\lv{6}
  \item fold a path \emph{product} with \(\prod_x w_x^{\#\text{partners}(x)}\) and combine arms by
        a quotient of products; \hfill\lv{6}
  \item solve CF~833D and hand-check it against the official sample. \hfill\lv{6}
\end{itemize}
\end{topicintro}

% ---------------------------------------------------------------
\subsection{The template, with the plugs pulled out}
% ---------------------------------------------------------------

\begin{definition}{Centroid divide-and-conquer, in general}{c18-beyond}
Fix, for the whole algorithm: a per-node \emph{feature map} \(x\mapsto \phi_c(x)\)
computed at centroid \(c\) that is \emph{additive} in the sense
\(\phi_c(\text{path } x\text{--}y)=\phi_c(x)+\phi_c(y)\) for a pair straddling \(c\); an
\emph{acceptance predicate} on \(\phi_c(x)+\phi_c(y)\) that is a conjunction of
inequalities each monotone in one coordinate; and a commutative \emph{accumulator}
\(\oplus\) over per-pair values with identity \(e\). Then
\(\mathit{solve}(S)\): choose centroid \(c\); gather each arm's node features; combine
arms through an ordered structure to fold \(\oplus\) over accepted straddling pairs
\emph{minus} the pairs inside a single arm; recurse on the pieces, multiplying the
answers. \(\bigO{\log n}\) levels make the total the per-centroid combine times
\(\log n\).
\end{definition}

\noindent\textbf{In plain words.} Find the balance node, describe a path by numbers that
\emph{add} when you glue its two halves at the centroid, keep the pairs whose summed
numbers pass a "both coordinates non-negative" test, and fold whatever you are asked to
fold. Everything in \(\secref{c18-count}\) and \(\secref{c18-index}\) was one instance;
CF~833D is the instance where the folded value is a product and the feature is a colour
imbalance.

\begin{insightbox}[Why a product still permits the same-arm subtraction]
For a count, \(\#\text{cross}=\#(\text{all})-\sum\#(\text{arm})\) because the qualifying
pairs split into "inside one arm" and "straddling." For a \emph{product} the pairs
partition the same way, so the product over the union is the product of the parts:
\(G(\text{all})=G_{\text{cross}}\cdot\prod_{\text{arm}}G(\text{arm})\). Divide out the arm
terms and you have the straddling product,
\(G_{\text{cross}}=G(\text{all})\cdot\prod_{\text{arm}}G(\text{arm})^{-1}\), where the
inverse is modular (Fermat, since the modulus \(10^9{+}7\) is prime). Subtraction becomes
division --- the only change a multiplicative accumulator forces.
\end{insightbox}

\begin{lstlisting}[caption={The reusable centroid divide-and-conquer skeleton with its three plug-ins named. Feature (additive per node), accept (2-D dominance), and the accumulator fold are the only things that change between 24.2, 24.3 and 24.4. ASCII only}, label={lst:c18-beyond}]
// template -- instantiate FEATURE / ACCEPT / FOLD for each problem.
void solve_comp(int entry){
    int tot = dfs_size(entry,-1);
    int c   = find_centroid(entry,-1,tot);        // lst:c18-build
    // 1. FEATURE: walk each arm, compute phi_c(x) for every node.
    vector<Feat> all = {(0,0, ELEM_IDENTITY)};    // the centroid: neutral feature
    for (v in adj[c]) if (!dead[v]) {
        Feat cur = feature_from(c, v);            // additive along c->v
        collect(cur, v, c);                        // DFS, skipping dead[]
    }
    // 2. ACCEPT: keep pair {x,y} iff phi(x)+phi(y) passes a dominance test.
    // 3. FOLD:    acc = fold over accepted cross-arm pairs;  subtract by the
    //             same-arm trick -- counts: acc -= fold(arm);  products: acc /= fold(arm).
    result_c = combine(all) MINUS_or_DIVIDE sum over arms of combine(arm);
    dead[c] = true;                               // wall off c, then recurse
    for (v in adj[c]) if (!dead[v]) solve_comp(v);
    combine_into_global(result_c);                // '+' for counts, '*' for products
}
// 24.2: Feat = depth,            ACCEPT = depth sum == K,   FOLD = count,  subtract.
// 24.3: Feat = (edges, weight),  ACCEPT = both <= budget,    FOLD = count,  subtract.
// 24.4: Feat = (2r-b, 2b-r, w),  ACCEPT = both >= 0,         FOLD = product, divide.
\end{lstlisting}

\begin{observation}[Reading a new statement as an instance]
Before reaching for code, ask the three plug-in questions. \emph{Feature:} is the path
property a \(\sum\) of per-edge contributions that can be regrouped as (arm1)+(arm2)?
\emph{Accept:} is the filter a monotone test on those sums (a threshold, a congruence, a
two-sided bound)? \emph{Fold:} is the answer count / sum / product / max over accepted
pairs? If any answer is "no," the centroid bucket will not directly carry it --- go to
\(\secref{c18-boundary}\) and pick a different engine.
\end{observation}

% ---------------------------------------------------------------
\subsection{Problem Lab: CF 833D -- a colour-balance product over all paths}
% ---------------------------------------------------------------

\solvedproblem{c18-833d}{Red-Black Cobweb}{CF 833D}
  {https://codeforces.com/problemset/problem/833/D}{CF 2800 L5}
  {centroid decomposition; colour imbalance as two additive dominance features; product accumulator folded as prod of w to the partner-count; same-arm subtraction by modular division}

\begin{problemstatement}[Red-Black Cobweb (statement and samples verified against the live problem)]
A tree on \(n\) nodes (\(2\le n\le 10^5\)); each edge has a \emph{clamminess}
\(x\) (\(1\le x\le 10^9{+}6\)) and a colour \(c\in\{0,1\}\) (0 = red, 1 = black). For an
unordered pair of nodes, look at its unique path: let \(r,b\) be the counts of red and
black edges on it. The pair is \emph{good} if the two counts "differ at most twice,"
i.e.\ \(\max(r,b)\le 2\min(r,b)\). The \emph{jelliness} is the product, over all good
pairs, of the product of clamminess along that path, modulo \(10^9{+}7\) (print \(1\) if
no good pair exists).

\textbf{Examples (official).}
\begin{verbatim}
Input:            Output:
5                 1265625
1 2 9 0
2 3 5 1
2 4 5 0
2 5 5 1

Input:            Output:
8                 452841614
1 2 7 1
2 3 4 1
3 4 19 1
5 1 2 0
6 2 3 0
7 3 3 0
8 4 4 0
\end{verbatim}
\end{problemstatement}

\begin{attemptgate}[Attempt it yourself first -- \lv{6}, give yourself 50 minutes]
Translate "differ at most twice" into algebra before touching the tree.
\peek{if \(r\) reds and \(b\) blacks, what two inequalities say \(\max(r,b)\le 2\min(r,b)\)?}
\nudge{\(r\le 2b\) and \(b\le 2r\), i.e.\ \(2r-b\ge 0\) and \(2b-r\ge 0\). Both left sides are
\emph{sums} of a per-edge \(+2/-1\) value, so they add along a path split at a centroid.}
\fullhint{Give each node at centroid \(c\) the feature \((u,v,w)\) where \(u=2r-b\),
\(v=2b-r\) (\(r,b\) from \(c\) to the node) and \(w\) is the clamminess product from \(c\)
to the node. A straddling pair \(\{x,y\}\) is good iff \(u_x+u_y\ge 0\) and \(v_x+v_y\ge 0\)
--- a 2-D dominance count, swept like \(\coderef{c18-293e}\). The path product is
\(w_x w_y\), so the fold is \(G=\prod_x w_x^{\#\text{partners}(x)}\). Compute \(G(\text{all})\)
and divide by \(\prod_{\text{arm}}G(\text{arm})\) (modular inverse); \(c\) enters as
\((0,0,1)\). Multiply \(G\) into the answer, recurse. \(\bigO{n\log^2 n}\).}
\end{attemptgate}

\solstep{Understanding}
The filter is not a distance bound; it is a \emph{ratio} bound between two colour counts.
Write it as the pair of linear inequalities \(2r\ge b\) and \(2b\ge r\). Define the
imbalance features \(u=2r-b\) and \(v=2b-r\). Because \(r\) and \(b\) count edges of the
path, and edges of a straddling path split into the two arms at \(c\), both \(u\) and
\(v\) are additive: \(u(x\text{--}y)=u_x+u_y\). So "good" is exactly
"\(u_x+u_y\ge0\) and \(v_x+v_y\ge0\)." The \(w\) (clamminess product) is multiplicative,
but \(\log\)-wise additive; we keep it as a plain product.

\solstep{Planning}
Per centroid: list every node of the component as \((u,v,w)\) (plus the centroid as
\((0,0,1)\)). Count, for each node \(x\), its partners \(y\ne x\) in \emph{any} arm with
\(u_x+u_y\ge0,v_x+v_y\ge0\) --- a dominance query (sort by \(u\), Fenwick over compressed
\(v\)). Fold \(G=\prod_x w_x^{\mathrm{cnt}_x}\). Do this for the whole component and for
each single arm, then combine as \(G(\text{all})/\prod_{\text{arm}}G(\text{arm})\) to keep
only cross-arm (and with-centroid) pairs.

\solstep{Implementing}
The dominance sweep returns \(\mathrm{cnt}_x\) per node, but \(x\) counts \emph{itself}
when it is inserted with \(u_x\ge0,v_x\ge0\); subtract that self-match (it is not a pair).
Modular exponentiation turns a partner count into the product \(w_x^{\mathrm{cnt}_x}\).

\begin{lstlisting}[caption={CF 833D: centroid decomposition. Each node carries the additive colour-imbalance features u=2r-b, v=2b-r and the clamminess product w. A 2-D dominance sweep counts partners; the accumulator is a product, combined across arms by modular division. dfs size and find centroid match lst:c18-build. ASCII only}, label={lst:c18-833d}]
#include <bits/stdc++.h>
using namespace std;
typedef long long ll;
const int MAXN = 100005;
const ll MOD = 1000000007LL;
int n;
struct Edge{ int to, col; ll x; };
vector<Edge> adj[MAXN];
bool dead[MAXN]; int sz[MAXN];
ll ans = 1;

ll modpow(ll b, ll e){ ll r=1%MOD; b%=MOD;
    while(e){ if(e&1) r=r*b%MOD; b=b*b%MOD; e>>=1; } return r; }

struct Feat{ ll u, v, w; };            // u=2r-b, v=2b-r, w=product along c->node

// G = product over accepted unordered pairs {x,y} of (w_x * w_y)
ll merge(const vector<Feat>& pts){
    int m = pts.size(); if(m <= 1) return 1;
    vector<ll> vs; for(auto&t:pts) vs.push_back(t.v);
    sort(vs.begin(),vs.end()); vs.erase(unique(vs.begin(),vs.end()),vs.end());
    int B = vs.size(); vector<int> bit(B+1,0);
    auto add=[&](int i,int d){ for(;i<=B;i+=i&-i) bit[i]+=d; };
    auto qry=[&](int i){ int s=0; for(;i>0;i-=i&-i) s+=bit[i]; return s; };
    vector<int> ins(m), ord(m); for(int i=0;i<m;i++) ins[i]=ord[i]=i;
    sort(ins.begin(),ins.end(),[&](int a,int b){return pts[a].u > pts[b].u;});   // by u desc
    sort(ord.begin(),ord.end(),[&](int a,int b){return pts[a].u < pts[b].u;});   // queries: T=-u desc
    ll res=1; int p=0;
    for(int qi=0; qi<m; qi++){
        int x=ord[qi]; ll T=-pts[x].u;                      // partners need u_y >= T
        while(p<m && pts[ins[p]].u >= T){                  // insert eligible u_y
            int bi=lower_bound(vs.begin(),vs.end(),pts[ins[p]].v)-vs.begin()+1;
            add(bi,1); p++;
        }
        int below = lower_bound(vs.begin(),vs.end(),-pts[x].v)-vs.begin(); // distinct v < -v_x
        ll cnt = p - qry(below);                           // inserted with v_y >= -v_x
        if(pts[x].u >= 0 && pts[x].v >= 0) cnt--;          // drop the self-match
        res = res * modpow(pts[x].w, cnt) % MOD;           // fold w_x^{cnt_x}
    }
    return res;
}
void collect(int u,int p,int r,int b,ll w,vector<Feat>& arm,vector<Feat>& all){
    arm.push_back({2*r-b, 2*b-r, w}); all.push_back({2*r-b, 2*b-r, w});
    for(auto&e:adj[u]) if(e.to!=p && !dead[e.to])
        collect(e.to,u, r+(e.col==0), b+(e.col==1), w*e.x%MOD, arm, all);
}
void process(int c){
    vector<Feat> all; all.push_back({0,0,1});               // centroid itself
    vector<vector<Feat>> arms;
    for(auto&e:adj[c]) if(!dead[e.to]){
        vector<Feat> arm;
        collect(e.to,c, e.col==0, e.col==1, e.x%MOD, arm, all);
        arms.push_back(arm);
    }
    ll tot = merge(all);
    for(auto&arm:arms) tot = tot * modpow(merge(arm), MOD-2) % MOD;   // divide same-arm
    ans = ans * tot % MOD;
}
int dfs_size(int u,int p){ sz[u]=1;
    for(auto&e:adj[u]) if(e.to!=p && !dead[e.to]) sz[u]+=dfs_size(e.to,u);
    return sz[u]; }
int find_centroid(int u,int p,int tot){
    bool moved=true;
    while(moved){ moved=false;
        for(auto&e:adj[u]) if(e.to!=p && !dead[e.to] && sz[e.to]>tot/2){ p=u;u=e.to;moved=true;break; } }
    return u; }
void decompose(int entry){
    int tot=dfs_size(entry,-1); int c=find_centroid(entry,-1,tot);
    process(c);
    dead[c]=true;
    for(auto&e:adj[c]) if(!dead[e.to]) decompose(e.to);
}
int main(){
    scanf("%d",&n);
    for(int i=1;i<n;i++){ int u,v,c; ll x; scanf("%d %d %lld %d",&u,&v,&x,&c);
        adj[u].push_back({v,c,x}); adj[v].push_back({u,c,x}); }
    decompose(1);
    printf("%lld\n", ans);
}
\end{lstlisting}

\begin{tracebox}[Sample 1 by hand -- centroid \(2\), features \((u,v,w)\) with \(u=2r-b,\ v=2b-r\)]
Edges: \(1\!-\!2\) red \(x9\); \(2\!-\!3\) black \(x5\); \(2\!-\!4\) red \(x5\);
\(2\!-\!5\) black \(x5\). At centroid \(2\) every other node is its own arm. Features:
\(c=2(0,0,1)\); \(1:(r1,b0)\Rightarrow(2,-1,9)\); \(3:(0,1)\Rightarrow(-1,2,5)\);
\(4:(1,0)\Rightarrow(2,-1,5)\); \(5:(0,1)\Rightarrow(-1,2,5)\). A straddling pair is good
iff \(u\)-sum \(\ge0\) and \(v\)-sum \(\ge0\):
\((1,3):2{-}1{=}1,\ {-}1{+}2{=}1\) ok; \((1,4):4,\ {-}2\) \(\times\); \((1,5):1,1\) ok;
\((3,4):1,1\) ok; \((3,5):{-}2,\ \times\); \((4,5):1,1\) ok; centroid pairs
\((2,\cdot)\) all fail (a one-colour path has \(2b\ge r\) or \(2r\ge b\) violated). Good
pairs: \(\{1,3\},\{1,5\},\{3,4\},\{4,5\}\) with path products \(45,45,25,25\). Partner
counts \(1{:}2,\,3{:}2,\,4{:}2,\,5{:}2,\,2{:}0\), so
\(G=9^2\!\cdot5^2\!\cdot5^2\!\cdot5^2 = 81\cdot625\cdot\ldots\). Grouping by pair
\(\prod w_x w_y=45\!\cdot\!45\!\cdot\!25\!\cdot\!25=1265625\). Each arm is a single node,
\(G(\text{arm})=1\). Recursion adds nothing. Answer \(1265625\). \(\checkmark\)
\end{tracebox}

\begin{complexitybox}[The product fold]
\(\mathit{merge}\) is a sort and a dominance sweep, \(\bigO{s\log s}\) on a size-\(s\)
list; it runs on \(\mathit{all}\) once and on each arm, and each modular inverse is
\(\bigO{\log \text{MOD}}\). A node lies in \(\bigO{\log n}\) components
(\(\thmref{c18-height}\)), so the total is \(\bigO{n\log^2 n + n\log n\log\text{MOD}}\),
comfortably inside 6 s for \(n=10^5\). The extra \(\log\) versus \(\secref{c18-count}\) is
the dominance structure the two-feature filter demands; the extra modular inverses are the
multiplicative accumulator's price.
\end{complexitybox}

\begin{warningbox}[What breaks the product version]
\begin{itemize}[nosep]
  \item \textbf{Dividing by zero.} \(G(\text{arm})\) is a product of \(w\)-powers; a
        single-node arm gives \(G=1\) (identity), never \(0\), since clamminess
        \(x\ge1\) and the modulus is prime. If you let a \(w\equiv0\pmod{\text{MOD}}\)
        slip in, the inverse vanishes --- but \(1\le x\le 10^9{+}6<\text{MOD}\) forbids it.
  \item \textbf{The self-match.} The sweep inserts node \(x\) and lets it dominate itself;
        forgetting the \(\mathit{cnt}--\) for \(u_x,v_x\ge0\) inflates an exponent and
        squares a \(w_x\) that should appear once per genuine partner.
  \item \textbf{Ratio vs difference.} "Differ at most twice" is a \emph{ratio}
        (\(\max\le2\min\)), not \(|r-b|\le2\). Misread it as a difference and the features
        are wrong; the sample's single-edge paths (diff \(1\)) would then be counted, but
        they are \emph{not} good here.
\end{itemize}
\end{warningbox}

\begin{hardeningbox}[Other problems on the same three plugs]
\begin{itemize}[nosep]
  \item \textbf{Sum of path products, not product.} Swap the accumulator to \(+\) and fold
        \(w_x\cdot(\text{partner weight-sum})\); the dominance sweep needs a prefix of
        \emph{products} rather than counts in the Fenwick --- still one structure per bucket.
  \item \textbf{Congruence filters.} A path divisible-by-\(m\) test (\(\secref{c18-count}\)'s
        CF~715C) is an \emph{equality} on the feature, answered with a hash map instead of a
        Fenwick: the template is the same, only the ordered structure changes to unordered.
  \item \textbf{Farthest-with-property.} Max-accumulator instances keep the centroid claim
        and let the bucket be a single running maximum --- no dominance needed when the
        feature is one-dimensional, which is why those problems read as easier.
\end{itemize}
\end{hardeningbox}

\begin{recallbox}[24.4 in one breath]
\emph{Pull the plugs out.} A centroid bucket can carry an \emph{additive} per-node feature,
accept pairs by a \emph{dominance} test, and fold \emph{any} commutative accumulator --- even
a product, whose same-arm subtraction becomes modular \emph{division}
(\(\defref{c18-beyond}\)). CF~833D reduces "colours differ by a factor of two" to
\(2r-b\ge0,\ 2b-r\ge0\) and answers the jelliness in \(\bigO{n\log^2 n}\)
(\(\coderef{c18-833d}\)). You now own the full divide-and-conquer-on-trees template.
\end{recallbox}
